% ============================================================================
% Fichier  : partie1/P01_C02_histoire_investigation.tex
% Titre    : Histoire de l'Investigation Numérique
% Auteur   : MINKA MI NGUIDJOÏ Thierry Emmanuel
% Version  : 1.0 -- Juin 2026
% Statut   : RÉDIGÉ
%
% Sources fusionnées :
%   - 1_histoire_investigation.tex (version précédente)
%   - chap2.pdf (archéologie foucaldienne, périodisation, formules,
%                BTK, Stuxnet, WannaCry, SolarWinds, perspectives)
%   - 2_grandes_affaires.tex (affaires complémentaires)
%   - M0_Cartographie_Menace_PNC.docx (Afrique, Cameroun, AFRIPOL,
%                MoMo fraud, SIM swapping, données ANTIC)
% ============================================================================

\chapter{Histoire de l'Investigation Numérique}
\label{chap:histoire}

\epigraph{%
    « Plus vous regarderez loin dans le passé, plus vous verrez loin dans
    le futur. »%
}{--- Winston Churchill, \textit{The Gathering Storm}, 1948}

\niveauChapitre{Fondamental}{%
    Ce chapitre prolonge le Chapitre~\ref{chap:phi-fond}~: la notion de
    régime de vérité numérique ($\vec{R}_t$) y est appliquée cas par cas.
    Aucun prérequis technique.%
}

\begin{boxobjectifs}
\begin{itemize}
    \item Reconstituer la \termecle{double hélice} de l'investigation
          numérique~: l'enroulement historique du technique et du juridique.
    \item Appliquer le modèle des régimes de vérité $\vec{R}_t$ à chaque
          période pour comprendre pourquoi les affaires ont produit les
          innovations qu'elles ont produites.
    \item Identifier les jalons africains et camerounais de l'histoire
          de la discipline~; positionner le Cameroun dans ce récit.
    \item Extraire de chaque affaire une leçon opérationnelle directement
          applicable à la pratique contemporaine.
    \item Anticiper les configurations futures des régimes de vérité à
          partir de la loi d'accélération épistémique.
\end{itemize}
\end{boxobjectifs}

\bigskip

% ============================================================================
\section*{La Double Hélice de l'Investigation Numérique}
% ============================================================================

L'histoire de l'investigation numérique ne se raconte pas en ligne droite.
Elle s'enroule selon une \termecle{double hélice}\index{double hélice} --- image
empruntée à la biologie moléculaire --- où deux brins s'entrelacent sans
jamais se confondre~: le brin \textbf{technique} (ce que les outils permettent
de faire) et le brin \textbf{juridique} (ce que le droit autorise à faire
et ce qu'il exige pour que ce soit recevable). Chaque rupture dans l'un des
brins provoque une tension dans l'autre. C'est cette tension répétée qui fait
avancer la discipline \cite{chap2pdf}.

\medskip

L'histoire se découpe en cinq régimes de vérité successifs, chacun déclenché
par un incident fondateur qui a contraint le régime précédent à se reconfigurer.
Le vecteur de dominance $\vec{R}_t$ de chaque époque, rappelé dans le tableau
synthétique ci-après, sera illustré affaire par affaire.

\begin{tableaucours}{p{2.2cm}p{2.5cm}p{3.0cm}p{3.8cm}}{%
    \tableauHeader{Période} &
    \tableauHeader{Incident fondateur} &
    \tableauHeader{Innovation} &
    \tableauHeader{Vecteur $\vec{R}_t$ dominant}
}
1970--1990 & Creeper / 414s & Trace électronique & $\alpha_T \approx 0.7$ \\
1990--2000 & Op. Sundevil / Mitnick & Chain of custody & $\alpha_J \approx 0.5$ \\
2000--2010 & Enron / McKinnon & Standards ISO/NIST & $\alpha_P \approx 0.5$ \\
2010--2020 & Silk Road / Panama & Analyse massive & $\alpha_T \approx 0.6$ \\
2020--\ldots & SolarWinds / IA & Preuves ZK / IA & $\alpha_T + \alpha_J \approx 0.8$ \\
\end{tableaucours}
\vspace{-0.8em}
\begin{center}{\small\itshape Tableau~2.1 --- La double hélice en cinq régimes
(d'après \cite{chap2pdf})}\end{center}

% ============================================================================
\section{Les Prémices (1970--1990)~: L'Émergence de la Trace Électronique}
% ============================================================================

\subsection{Contexte socio-technique du premier régime}

Cette période correspond à ce que Manovich appelle « la logique culturelle
de l'ordinateur personnel »~: le numérique est le territoire d'une élite
technique restreinte, les crimes sont commis par curiosité autant que par
appât du gain, et le cadre juridique est absent ou analogique. Le régime
de vérité est dominé par la technique ($\alpha_T \approx 0.7$)~: l'expert
technique isolé est l'unique autorité épistémique, et ses méthodes ne sont
soumises à aucune contrainte procédurale.

\medskip

Les premières investigations sont de l'\textbf{artisanat}~: chaque analyste
invente ses propres méthodes, documente selon son intuition, conclut selon
son expertise. La preuve paradigmatique de l'époque est le log système~---
une trace brute, technique, dont la signification est accessible seulement
aux initiés.

\begin{tableaucours}{p{2.8cm}p{8.2cm}}{%
    \tableauHeader{Dimension} &
    \tableauHeader{Caractéristiques du régime 1970--1990}
}
Technologie & Mainframes, premiers PC, réseau ARPANET, absence de sécurité \\
Juridique & Absence de cadre spécifique~; application analogique du droit pénal existant \\
Social & Fracture numérique massive~; les hackers sont des héros ou des démons \\
Pratique & Méthodes ad hoc~; expertise individuelle non transmissible~; aucun standard \\
\end{tableaucours}
\vspace{-0.8em}
\begin{center}{\small\itshape Tableau~2.2 --- Contexte socio-technique des prémices
(d'après \cite{chap2pdf})}\end{center}

\subsection{The Creeper (1971)~: le premier artefact forensique}

\textbf{The Creeper} n'est pas un crime~: c'est une expérience. Bob Thomas,
ingénieur chez BBN Technologies, écrit en 1971 un programme capable de
se déplacer d'ordinateur en ordinateur sur ARPANET en affichant le message
\texttt{I'M THE CREEPER: CATCH ME IF YOU CAN!}. La réponse --- \textbf{The
Reaper}, programme écrit pour supprimer The Creeper --- constitue la
première \emph{investigation active} de l'histoire numérique~: pour
attraper The Creeper, il fallait comprendre son comportement, sa propagation,
ses traces.

\medskip

La leçon n'est pas technique~; elle est épistémologique~: \textbf{tout
programme laisse des traces}, et ces traces peuvent être lues par qui sait
les interpréter. Ce principe, évident aujourd'hui, était révolutionnaire
en 1971.

\subsection{L'affaire des 414s (1983)~: premier paradigme investigatif}

En 1983, six adolescents de Milwaukee, surnommés les \termecle{414s} d'après
leur indicatif régional, pénètrent dans 60~systèmes informatiques incluant le
Los Alamos National Laboratory et le Memorial Sloan-Kettering Cancer Center.
Ils sont arrêtés, non par des techniques forensiques sophistiquées, mais parce
que l'un d'eux se vante devant ses pairs --- un informateur prévient la police.

\medskip

Du point de vue forensique, cette affaire établit que les traces réseau
constituent une preuve légitime, créant le premier régime de vérité basé
sur l'adresse IP \cite{chap2pdf}~:

\begin{equation}
P_{\text{réseau}} = \bigl\{\text{logs},\; \text{adresses IP},\;
\text{horodatages}\bigr\} \in \mathcal{P}_{\text{légitime}}
\label{eq:preuve-reseau}
\end{equation}

\begin{boxCRO}
\textbf{Analyse CRO --- Affaire des 414s}

L'investigation repose sur l'exploitation des logs système~: aucune protection
de la vie privée n'est invoquée. \cro{$\downarrow$}{$\uparrow$}{$\downarrow$} ---
Les preuves sont techniquement solides mais le cadre juridique n'existe pas encore
pour les qualifier. L'absence de \loicam{} ou de \textit{Computer Fraud and Abuse
Act} avant 1986 illustre le retard du brin juridique sur le brin technique.
\end{boxCRO}

\textbf{Impact~:} création du \textit{Computer Fraud and Abuse Act} aux
États-Unis en 1986. Prise de conscience politique mondiale de la nécessité
d'un cadre juridique propre aux crimes informatiques.

\subsection{L'affaire BTK Killer (2005, racines 1974)~: les métadonnées comme preuve}

Dennis Rader, tueur en série surnommé BTK (\textit{Bind, Torture, Kill}),
avait communiqué avec la police pendant 30~ans sans être identifié. En
janvier 2005, il envoie une disquette contenant un fichier
\texttt{Test.A.rtf}. Les enquêteurs du FBI extraient les
\termecle{métadonnées}\index{métadonnées} du fichier avec un outil basique~:
le champ « Auteur » contient \texttt{Dennis}, et la propriété « Dernière
organisation » indique \texttt{Christ Lutheran Church}. En croisant avec
les membres actifs de cette paroisse, Rader est identifié en quelques
heures.

\medskip

Cette affaire illustre un principe cardinal~: \textbf{la valeur probatoire
d'un artefact numérique peut dépasser celle de son contenu}. Le message
de Rader n'incriminait personne~; ses métadonnées l'ont trahi \cite{chap2pdf}~:

\begin{equation}
V_{\text{probatoire}}(\text{métadonnées}) \;\geq\;
V_{\text{probatoire}}(\text{contenu})
\label{eq:meta-superieur}
\end{equation}

\begin{boxartefact}{Métadonnées de fichier Office}
\textbf{Localisation~:} \texttt{Propriétés} du fichier $\to$
\texttt{Détails} (Windows) ou extraction via ExifTool/Autopsy.

\textbf{Informations typiquement contenues~:} auteur, organisation, date de
création, date de dernière modification, logiciel utilisé, parfois chemin
du fichier sur l'ordinateur source, nom d'imprimante.

\textbf{Commande d'extraction~:}
\lstinline[language=bash_forensique]{exiftool fichier.docx}

\textbf{Vulnérabilité~:} la plupart des utilisateurs ignorent que ces
données existent et n'ont pas l'habitude de les supprimer avant d'envoyer
un fichier. L'investigateur doit systématiquement les extraire.
\end{boxartefact}

% ============================================================================
\section{La Professionnalisation (1990--2000)~: Naissance d'une Discipline}
% ============================================================================

\subsection{La révolution Internet et ses conséquences épistémiques}

L'émergence du web transforme radicalement le régime de vérité. Le brin
juridique ($\alpha_J$) et le brin pratique ($\alpha_P$) deviennent dominants~:
l'investigation ne peut plus être artisanale quand elle touche des millions
d'internautes. La \textbf{professionnalisation} au sens d'Abbott correspond
à ce moment où un groupe construit sa juridiction professionnelle autour~:
d'un corps de connaissances formalisé~; de méthodes standardisées et
reproductibles~; d'une reconnaissance juridique et institutionnelle~; et
d'un code déontologique contraignant \cite{chap2pdf}.

\subsection{L'Opération Sundevil (1990)~: le choc de la standardisation}

Le 8~mai~1990, le Secret Service américain lance la plus grande opération
contre la cybercriminalité jamais conduite~: saisie de 42~systèmes
informatiques dans 14~villes simultanément. L'opération est un succès
opérationnel et un désastre juridique. La majorité des poursuites
s'effondrent~: les preuves saisies n'ont pas été préservées selon des
protocoles cohérents~; les hash n'existent pas encore comme concept
forensique~; plusieurs systèmes saisis appartiennent à des innocents ciblés
par erreur.

\medskip

L'Opération Sundevil représente, selon Durkheim, un \textbf{fait social total}~:
elle cristallise tous les enjeux de la nouvelle société numérique et révèle
l'impossibilité de continuer avec des méthodes artisanales. Elle conduit
directement à la création de l'\sigle{IOCE} (\textit{International Organization
on Computer Evidence}) en 1995 --- première institution dédiée à la
standardisation de la forensique numérique.

\begin{boxjuris}
\textbf{Héritage normatif de l'Opération Sundevil~:}
\begin{itemize}
    \item Création de l'\sigle{IOCE} (1995) --- premiers principes de collecte
          de preuves numériques.
    \item Fondation du \sigle{DFRWS} (\textit{Digital Forensics Research
          Workshop}) (2001) --- qui produira le premier modèle de processus
          forensique.
    \item Inspiration directe pour la Convention de Budapest (2001).
\end{itemize}
\end{boxjuris}

\subsection{L'arrestation de Kevin Mitnick (1995)~: institution de la chaîne de custody}

Le 15~février~1995, Kevin Mitnick --- « le hacker le plus recherché du monde »
--- est arrêté à Raleigh (Caroline du Nord) grâce au travail de Tsutomu
Shimomura. L'affaire dure depuis 1993~: Mitnick a utilisé le \textit{IP
spoofing}, le \textit{TCP session hijacking} et l'ingénierie sociale pour
pénétrer des dizaines de systèmes militaires et corporatifs.

\medskip

La contribution épistémologique de cette affaire est plus importante que
sa dimension médiatique~: elle institue la \textit{chain of custody}
comme \textbf{condition nécessaire de la validité épistémique} d'une
preuve numérique \cite{chap2pdf}~:

\begin{equation}
\mathrm{Validit}(P) \;\Leftrightarrow\;
\exists\, \mathcal{C} \;\text{ telle que }\; \int \mathrm{grit}(\mathcal{C})\,dt = 1
\label{eq:validite-custody}
\end{equation}

où $\mathcal{C}$ représente la chaîne de custody et $\mathrm{grit}(\mathcal{C})$
sa densité documentaire. Informellement~: une preuve n'est valide que si elle
est accompagnée d'une documentation continue de sa traçabilité.

\begin{boxCRO}
\textbf{Analyse CRO --- Arrestation Mitnick}

L'affaire Mitnick mobilise les trois dimensions sous forte tension.
$\mathcal{C}$~: les écoutes téléphoniques étendues soulèvent des questions
de vie privée qui alimenteront la littérature juridique pendant une décennie.
$\mathcal{R}$~: l'innovation de Shimomura (analyse de honeypot, corrélation
temporelle avancée) est robuste et reproductible. $\mathcal{O}$~: la chaîne
de custody, bien que perfectible, est suffisamment documentée pour permettre
une condamnation.
\cro{$\downarrow$ tendu}{$\uparrow$}{$\uparrow$ émergent}
\end{boxCRO}

% ============================================================================
\section{La Standardisation (2000--2010)~: Rationalisation Industrielle}
% ============================================================================

\subsection{Le tournant du millénaire et les standards fondateurs}

Cette période correspond à ce que Beck appelle « la société du risque »~:
le numérique est devenu infrastructure critique, et les États comme les
entreprises comprennent que les méthodes artisanales sont un risque systémique.
Quatre standards fondateurs émergent simultanément, créant une infrastructure
normative sans précédent~:

\begin{tableaucours}{p{2.4cm}p{2.2cm}p{3.5cm}p{3.2cm}}{%
    \tableauHeader{Standard} &
    \tableauHeader{Portée} &
    \tableauHeader{Contribution principale} &
    \tableauHeader{Acteurs moteurs}
}
\norme{ISO~27037} & International &
    Méthodologie unifiée de collecte & Organisations internationales \\
\norme{NIST SP 800-86} & États-Unis &
    Gouvernance processuelle complète & Agences gouvernementales \\
\norme{RFC~3227} & Internet &
    Ordre de volatilité --- bonnes pratiques & Communauté technique \\
\norme{ACPO Guide} & Royaume-Uni &
    4~principes déontologiques & Forces de l'ordre britanniques \\
\end{tableaucours}
\vspace{-0.8em}
\begin{center}{\small\itshape Tableau~2.3 --- Standards fondateurs de l'ère 2000--2010
(d'après \cite{chap2pdf})}\end{center}

\subsection{L'affaire Enron (2001)~: naissance de l'\textit{e-discovery}}

La faillite d'Enron produit l'investigation numérique la plus massive de
l'histoire jusqu'alors~: 500~000~documents électroniques doivent être analysés,
croisés, et présentés en justice. Les méthodes manuelles sont impossibles~;
la réponse est la création des premiers outils d'\termecle{analyse automatisée}
(précurseurs du \textit{Technology-Assisted Review}). L'affaire établit que
les résultats d'un algorithme d'analyse constituent des preuves légitimes
si l'algorithme est validé et documenté~:

\begin{equation}
P_{\text{algo}} = \mathcal{A}(\mathcal{D}) \in \mathcal{P}_{\text{légitime}}
\quad \text{si} \quad \mathrm{si}(\mathcal{A}) > \theta
\label{eq:preuve-algo}
\end{equation}

où $\mathrm{si}(\mathcal{A})$ est l'indice de scientificité de l'algorithme et
$\theta$ le seuil d'acceptabilité judiciaire. Les amendements aux
\textit{Federal Rules of Civil Procedure} (2006) codifient ces principes
en droit positif américain, ouvrant la voie à l'automatisation de l'investigation.

\subsection{L'affaire Gary McKinnon (2002--2012)~: souveraineté numérique}

Gary McKinnon, informaticien britannique autiste, passe deux ans à infiltrer
97~serveurs militaires américains à la recherche de preuves d'existence
extraterrestre. L'affaire dure dix~ans, non pour des raisons techniques (McKinnon
est rapidement identifié), mais parce qu'elle cristallise un conflit de
\textbf{souveraineté numérique}~: les États-Unis réclament l'extradition
pour 70~ans de prison potentiels~; le Royaume-Uni refuse pendant une décennie
au nom du droit à la santé de McKinnon.

\medskip

La leçon forensique est la corrélation multi-juridictionnelle~: quand une
attaque traverse des frontières, les logs se trouvent dans des pays différents,
les règles d'admissibilité divergent, et la coopération judiciaire internationale
devient le vrai défi --- pas la technique.

% ============================================================================
\section{Le Big Data et le Cloud (2010--2020)~: L'Échelle Critique}
% ============================================================================

\subsection{Le changement d'échelle qualitatif}

Cette période correspond à ce que Boyd et Crawford appellent « l'ère des
big data » où l'échelle change la nature même de l'investigation. Au-delà
d'un certain seuil $N_{\text{critique}}$, l'analyse ne peut plus être humaine~:

\begin{equation}
\lim_{N \to \infty}
\frac{\text{Analyse humaine}}{\text{Analyse algorithmique}} = 0
\label{eq:bigdata}
\end{equation}

Cette équation a une implication directe sur le Trilemme CRO~: quand
l'analyse est déléguée à un algorithme opaque, $\mathcal{O}$ (opposabilité)
est menacée --- l'avocat de la défense peut contester une conclusion
qu'aucun expert humain ne peut entièrement expliquer.

\subsection{L'affaire Silk Road (2013)~: forensique blockchain}

Ross Ulbricht exploite Silk Road, première grande place de marché du dark
web, depuis 2011. L'investigation du FBI combine trois innovations~: la
\textbf{forensique Tor} (corrélation de flux pour désanonymiser les serveurs),
la \textbf{forensique blockchain} (analyse du graphe de transactions Bitcoin
pour remonter les flux financiers), et l'\textbf{OSINT} (Ulbricht a publié
des messages identifiables sur un forum public). La saisie de 144~000~bitcoins
(environ 28~millions~USD à l'époque, bien plus aujourd'hui) établit que
les cryptomonnaies ne sont pas anonymes~: elles sont pseudonymes, et le
pseudonymat se brise par corrélation de métadonnées.

\begin{boiteRemarque}{La Blockchain comme registre forensique}
La blockchain Bitcoin est, du point de vue forensique, le registre comptable
le plus complet jamais créé~: chaque transaction est permanente, publique,
et cryptographiquement liée à toutes les précédentes. La forensique blockchain
exploite cette propriété en construisant le graphe complet des flux pour
identifier les points d'échange (exchanges) où les cryptomonnaies sont
converties en monnaie réelle, permettant l'identification des acteurs.
\end{boiteRemarque}

\subsection{Les Panama Papers (2016)~: investigation transfrontière}

La fuite de 2,6~To de données (11,5~millions de documents) du cabinet
Mossack Fonseca constitue la plus grande investigation de données de
l'histoire. 376~journalistes de 76~pays collaborent pendant un an
sur un corpus unifié. L'impact forensique est l'émergence de
l'\termecle{investigation transnationale}\index{investigation transnationale}
comme nouveau paradigme~:

\begin{equation}
I_{\text{global}} = \sum_{i=1}^{n} I_{\text{locale}}^{i} +
C_{\text{collaboration}} - L_{\text{juridique}}
\label{eq:investigation-globale}
\end{equation}

où $C_{\text{collaboration}}$ est le bénéfice de la coopération et
$L_{\text{juridique}}$ les frictions des barrières légales entre juridictions.
La leçon~: la coordination internationale compense les barrières juridiques,
mais cette coordination exige des protocoles techniques et déontologiques
communs --- exactement ce que ce cours enseigne.

% ============================================================================
\section{L'Ère Post-Quantique et IA (2020--\ldots)~: Le Grand Saut}
% ============================================================================

\subsection{La double révolution IA / quantique}

Nous entrons dans ce que l'on pourrait appeler l'ère de la « guerre
algorithmique »~: l'IA est simultanément l'arme de l'attaquant et
le bouclier de l'investigateur. Deux ruptures parallèles redéfinissent
les fondamentaux.

\medskip

La première rupture est l'\textbf{émergence du régime de vérité algorithmique}~:
la vérité numérique est désormais établie par des systèmes d'IA dont le
fonctionnement peut être incompréhensible pour les humains (\textit{black box
problem}). Cela crée une tension épistémique fondamentale~: peut-on fonder
une condamnation pénale sur un résultat algorithmique inexplicable~?

La seconde rupture est la \textbf{menace quantique sur la cryptographie}~:
l'algorithme de Shor rend théoriquement cassables RSA et ECC sur un ordinateur
quantique suffisant. Les preuves numériques collectées aujourd'hui et
reposant sur ces algorithmes pourraient être invalidées rétroactivement.

\subsection{L'attaque SolarWinds (2020)~: IA contre IA}

L'attaque SolarWinds est la première illustration documentée d'une
\textbf{attaque \textit{living-off-the-land}} à l'échelle des États~:
les attaquants (groupe APT29, attribué aux services de renseignement russes)
compromettent la chaîne d'approvisionnement logicielle de SolarWinds en
injectant un code malveillant dans une mise à jour légitime. 18~000~organisations
téléchargent la mise à jour infectée~; la compromission reste non détectée
pendant \textbf{9~mois} malgré des outils de sécurité sophistiqués.

\medskip

La leçon forensique centrale~: la détection classique basée sur des signatures
est morte. Seule l'\textbf{analyse comportementale basée sur l'IA} permet
de détecter ce type d'attaque --- et encore, après le fait. L'attribution
reste incertaine~: les techniques d'obfuscation masquent l'origine, et
l'investigateur doit composer avec une probabilité, pas une certitude.

\begin{boxCRO}
\textbf{Analyse CRO --- Attaque SolarWinds}

$\mathcal{C}$~: l'investigation exige l'accès aux systèmes internes de
18~000~organisations --- tension massive avec la vie privée des clients.
$\mathcal{R}$~: la reproductibilité est problématique~; les outils IA
utilisés pour la détection ne produisent pas des conclusions entièrement
auditables. $\mathcal{O}$~: l'attribution probabiliste basée sur l'IA
est encore mal acceptée par les juridictions pénales~; aucune condamnation
n'a suivi. \cro{$\downarrow$ critique}{$\sim$ partielle}{$\downarrow$ problématique}

Ce cas illustre la \textbf{crise d'opposabilité du régime algorithmique}~:
les outils sont puissants, les conclusions solides techniquement, mais leur
défendabilité devant un tribunal reste à construire.
\end{boxCRO}

\subsection{Affaire WannaCry (2017)~: immunité collective numérique}

Le~12~mai~2017, le ransomware WannaCry infecte 300~000~ordinateurs dans
150~pays en 72~heures, exploitant la vulnérabilité \textit{EternalBlue}
volée à la NSA. La réponse forensique de Marcus Hutchins --- un chercheur
de 22~ans qui identifie et active un \textit{kill switch} caché dans le code,
arrêtant la propagation --- illustre un principe nouveau~: la réponse aux
incidents majeurs devient une entreprise \textbf{décentralisée et collaborative}
similaire à l'immunité collective en épidémiologie~:

\begin{equation}
R_0 = \beta \cdot \gamma^{-1}
\label{eq:wannacry-R0}
\end{equation}

où $\beta$ est le taux d'infection et $\gamma$ le taux de désinfection.
L'analyse du \textit{Domain Generation Algorithm} (DGA) de WannaCry par
des dizaines d'équipes en parallèle constitue un modèle d'investigation
distribuée qui préfigure les pratiques actuelles du partage de renseignement
sur les menaces (\textit{Threat Intelligence}).

% ============================================================================
\section{L'Afrique et le Cameroun dans l'Histoire de la Discipline}
% ============================================================================

\subsection{Une histoire trop souvent absente des manuels}

Les histoires standard de l'investigation numérique sont presque
exclusivement nord-américaines et européennes. C'est une distorsion~:
l'Afrique n'est pas un spectateur de la cybercriminalité mondiale, elle
en est simultanément une cible croissante, un terrain d'opération, et
--- de plus en plus --- une source d'innovation forensique. Cette section
comble ce vide historiographique.

\subsection{Les fraudes 419 (1990s--présent)~: la plus longue affaire forensique africaine}

Les escroqueries dites \termecle{419}\index{419, fraude} --- du numéro de
l'article du Code pénal nigérian qui les prohibe --- constituent la première
grande catégorie de cybercriminalité d'origine africaine à portée mondiale.
Dès les années 1990, des réseaux nigérians envoient par email des millions de
messages proposant des fortunes fictives. La réponse investigative révèle les
limites fondamentales du régime de vérité de l'époque~: les suspects sont en
Afrique, les victimes en Europe et en Amérique, les preuves traversent des
dizaines de juridictions, et aucune convention de coopération n'existe.

\medskip

Cette affaire permanente a deux héritages directs pour notre cours~: elle a
forcé la Convention de Budapest (2001) à inclure une dimension coopération
internationale, et elle a conduit le Cameroun à adopter la
\loicam{2010/012}, qui criminalise explicitement ces pratiques en droit
positif camerounais.

\subsection{Le Cameroun et la Loi 2010/012~: entrée dans le régime de standardisation}

Le 21~décembre~2010, le Cameroun adopte la \textbf{Loi N°2010/012 relative
à la cybersécurité et la cybercriminalité}. Cette date marque l'entrée
formelle du Cameroun dans le régime de standardisation mondial, avec un
décalage de dix ans sur les pays pionniers --- mais avec l'avantage de
bénéficier du retour d'expérience d'une décennie d'application internationale.
La même année, l'\sigle{ANTIC} est opérationnalisée comme autorité compétente.

\begin{boxjuris}
\textbf{Jalons camerounais de l'histoire forensique~:}
\begin{itemize}
    \item \textbf{2010}~: Loi 2010/012 (cybersécurité et cybercriminalité)
          et Loi 2010/013 (communications électroniques).
    \item \textbf{2011}~: Création et montée en puissance de l'\sigle{ANTIC}
          comme autorité de régulation et de certification.
    \item \textbf{2016}~: Compétence du Tribunal Criminel Spécial (\sigle{TCS})
          étendue à certaines infractions cyber.
    \item \textbf{2017}~: Opérationnalisation du \sigle{CIRT} national
          (Computer Incident Response Team).
    \item \textbf{2019--présent}~: Multiplication des affaires de fraude
          Mobile Money (SIM swapping, BEC) --- première jurisprudence
          camerounaise en investigation numérique.
    \item \textbf{2024}~: Loi sur la protection des données personnelles ---
          nouveau cadre pour les investigations impliquant des données privées.
\end{itemize}
\end{boxjuris}

\subsection{L'Opération FALCON II (2022)~: la coopération africaine en action}

En novembre 2022, Interpol et Afripol coordonnent l'Opération FALCON II
impliquant 11~pays africains dont le Cameroun. Le résultat~: arrestation
de 11~individus, démantèlement du groupe \textbf{SilverTerrier} (Nigeria),
responsable de plus de 1,7~million d'emails frauduleux et de pertes évaluées
à plusieurs centaines de millions de dollars.

\medskip

Cette opération représente une rupture de paradigme pour l'Afrique~:
la coopération régionale forensique produit des résultats mesurables. Elle
prouve que le modèle de l'investigation transnationale --- identifié dans les
Panama Papers --- est applicable sur le continent, avec les outils et le
cadre normatif existants \cite{m0pnc}.

\begin{boxscenario}{Fraude Mobile Money --- Cameroun, 2021}
\textbf{Faits~:} Un réseau de SIM swapping opère dans la région du Centre.
Méthode~: collecte des données personnelles des victimes via de faux sondages
diffusés sur WhatsApp, puis substitution de la SIM dans une agence MTN
complice. Résultat~: 47~millions FCFA détournés en 3~semaines, 23~victimes.

\textbf{Arrestations~:} 4~personnes dont 2~employés MTN. Qualification
juridique~: escroquerie informatique (\loicam{2010/012}, art.~78)
et complicité.

\textbf{Leçon forensique~:} La preuve numérique centrale n'est pas le
message WhatsApp (contenu) mais les métadonnées de connexion (qui s'est
connecté au système MTN, depuis quelle adresse IP, à quelle heure) et
les logs de transactions Mobile Money. L'enquête illustre que la chaîne
de custody doit commencer chez l'opérateur télécom --- avant même la
saisie du téléphone du suspect.

\textbf{Analyse CRO~:} \cro{$\downarrow$}{$\uparrow$}{$\uparrow$} --- Les logs
opérateur sont techniquement fiables ($\mathcal{R}$ haute) et leur
communication est encadrée par l'art.~30 de la Loi~2010/012 (réquisition
judiciaire), garantissant $\mathcal{O}$. La vie privée des victimes
($\mathcal{C}$) est sacrifiée dans la limite du nécessaire, ce qui est
justifiable et proportionnel.
\end{boxscenario}

% ============================================================================
\section{Perspectives~: la Loi d'Accélération Épistémique}
% ============================================================================

\subsection{Le temps entre les régimes se raccourcit}

L'observation la plus frappante de l'histoire que nous venons de parcourir
est que le temps entre les ruptures de paradigme se réduit
exponentiellement \cite{chap2pdf}~:

\begin{equation}
\Delta t_{n+1} = k \cdot \Delta t_n \quad \text{avec} \quad 0 < k < 1
\label{eq:acceleration}
\end{equation}

Le passage du régime technique au régime juridique a pris vingt ans (1970--1990).
Du juridique au standardisé~: dix ans. Du standardisé au computationnel~: dix
ans. Du computationnel au post-quantique/IA~: probablement cinq ans. La
conséquence pratique est directe~: un investigateur formé uniquement sur les
outils actuels sera partiellement obsolète dans moins de cinq ans. La
compréhension des \emph{régimes} --- pas des outils --- est la seule
formation qui résiste à cette accélération.

\subsection{Trois scénarios pour 2030--2050}

L'analyse prospective identifie trois régimes plausibles \cite{chap2pdf}~:

\begin{enumerate}
    \item \textbf{Régime quantique dominant~:} le calcul quantique devient
          accessible. RSA et ECC sont cassées. Toute la forensique des preuves
          collectées avant 2030 doit être revue. Les protocoles ZK et les
          algorithmes post-quantiques standardisés par le NIST (CRYSTALS-Kyber,
          CRYSTALS-Dilithium) deviennent la nouvelle norme.

    \item \textbf{Régime IA investigatif~:} les LLM et les agents autonomes
          conduisent les premières phases de triage forensique sans intervention
          humaine. La question d'opposabilité devient centrale~: peut-on présenter
          en justice une preuve identifiée par un agent IA qu'aucun expert humain
          ne peut entièrement auditer~?

    \item \textbf{Régime post-humain~:} existe un point $t_s$ (\textit{singularité
          investigative}) où l'IA investigative dépasse durablement les capacités
          humaines. Ce régime pose la question fondamentale de la responsabilité~:
          si l'IA se trompe et conduit à une condamnation injuste, qui en est
          responsable~?
\end{enumerate}

\medskip

Ces scénarios ne sont pas des spéculations~: ils se préparent maintenant.
Les Parties~V et~VI de ce cours sont construites précisément pour équiper
l'étudiant face aux deux premiers scénarios.

\separateur

\begin{boxresume}
\textbf{Ce qu'il faut retenir de ce chapitre~:}

\begin{itemize}
    \item L'histoire de l'investigation numérique suit une
          \textbf{double hélice}~: le brin technique et le brin juridique
          s'enroulent et se contraignent mutuellement. Chaque rupture dans
          l'un provoque une réorganisation dans l'autre.

    \item Cinq régimes successifs~: \textbf{technique} (1970--1990),
          \textbf{juridique} (1990--2000), \textbf{standardisé} (2000--2010),
          \textbf{computationnel} (2010--2020), \textbf{post-quantique/IA}
          (2020--\ldots). Chacun est déclenché par un incident fondateur.

    \item Leçons opérationnelles permanentes~: les métadonnées peuvent
          dépasser le contenu en valeur probatoire (BTK)~; la chaîne de
          custody est une condition de validité, pas un formalisme
          (Mitnick)~; les algorithmes d'analyse sont des preuves à documenter
          (Enron)~; le pseudonymat blockchain se brise par corrélation
          (Silk Road)~; l'attribution probabiliste reste à la limite de
          l'opposabilité (SolarWinds).

    \item L'Afrique et le Cameroun ont leur propre histoire~: les fraudes~419,
          la \loicam{2010/012} (2010), le SIM swapping, l'Opération FALCON~II.
          Cette histoire n'est pas un appendice de l'histoire mondiale --- elle
          en est une dimension constitutive.

    \item La \textbf{loi d'accélération épistémique}
          ($\Delta t_{n+1} = k \cdot \Delta t_n$) prédit que les ruptures
          s'accélèrent. Maîtriser les régimes est plus durable que maîtriser
          les outils.
\end{itemize}
\end{boxresume}

\bigskip

\begin{boxexo}[title={\ding{45}\quad Exercice~2.1 --- Analyse d'un incident par son régime \niveaubase}]
Pour chacun des incidents suivants, identifiez le régime de vérité dominant,
le brin (technique ou juridique) qui a été le plus sollicité, et la leçon
forensique principale~:

\begin{enumerate}[(a)]
    \item Un journaliste publie un reportage sur la corruption d'un ministre
          en utilisant des emails obtenus par un lanceur d'alerte. Le tribunal
          accepte les emails comme preuves. Année~: 2004.
    \item Une entreprise camerounaise perd 80~millions FCFA via un BEC. La
          police obtient les logs de connexion auprès de l'opérateur internet
          sur réquisition judiciaire. Année~: 2022.
    \item Un analyste utilise un modèle LLM pour classer automatiquement
          50~000~messages WhatsApp dans une affaire de trafic de drogue.
          L'avocat de la défense conteste la méthode. Année~: 2025.
\end{enumerate}
\end{boxexo}

\begin{boxexo}[title={\ding{45}\quad Exercice~2.2 --- La double hélice appliquée \niveaumoyen}]
L'affaire Stuxnet (2010) est un cas où un programme malveillant, conçu par
des services étatiques, a physiquement endommagé des centrifugeuses iraniennes
d'enrichissement d'uranium. Son investigation a mobilisé simultanément~:
reverse engineering de code polymorphe (4~zero-days simultanés), attribution
géopolitique, droit international, et éthique de la guerre numérique.

\begin{enumerate}[(a)]
    \item Appliquez le modèle $\vec{R}_t$ à cette affaire~: quelles composantes
          $(\alpha_T, \alpha_J, \alpha_S, \alpha_P)$ sont les plus sollicitées~?
    \item Identifiez deux tensions CRO spécifiques à cette affaire et proposez
          comment un investigateur mandaté par une institution internationale
          devrait les gérer.
    \item En quoi cette affaire préfigure-t-elle le « régime IA contre IA »
          décrit en Section~2.5~?
\end{enumerate}
\end{boxexo}

\begin{boxexo}[title={\ding{45}\quad Exercice~2.3 --- Prospective \niveauexpert}]
La loi d'accélération épistémique ($\Delta t_{n+1} = k \cdot \Delta t_n$)
suggère que le prochain régime de vérité émergera vers 2027--2030.

\begin{enumerate}[(a)]
    \item Sur la base des tendances actuelles (IA générative, informatique
          quantique, blockchain forensique), construisez un scénario de rupture
          plausible~: quel incident fondateur pourrait déclencher le prochain
          changement de régime~?
    \item Quelles compétences un investigateur formé aujourd'hui devrait-il
          développer pour rester pertinent dans ce nouveau régime~?
    \item Positionnez le Cameroun dans ce scénario~: quels atouts et quelles
          vulnérabilités spécifiques le pays présente-t-il face à cette
          transition~?
\end{enumerate}
\end{boxexo}

\bigskip

\begin{boxinfo}
\textbf{Transition.} Le Chapitre~\ref{chap:grandes-affaires} approfondit
cinq affaires choisies pour leur valeur pédagogique maximale --- chacune
illustrant une innovation forensique qui a redéfini la pratique. Le
Chapitre~\ref{chap:fond-theo} formalisera ensuite les outils théoriques
(théorie de l'information, théorie des graphes, cryptographie) sur lesquels
repose la forensique moderne.
\end{boxinfo}
